%% ============================================================
%%  Chapter 2 — Background
%% ============================================================
\chapter{Background and Related Work}
\label{ch:background}

\section{Reinforcement Learning}

Reinforcement learning (RL) is a paradigm in which an \emph{agent} learns to
make decisions by interacting with an \emph{environment} and receiving scalar
\emph{rewards}~\citep{sutton2018reinforcement}.
The goal is to find a policy $\pi(a \mid s)$ --- a mapping from states to
actions --- that maximises the expected cumulative discounted return:
%
\begin{equation}
  G_t = \sum_{k=0}^{\infty} \gamma^k \cdot r_{t+k+1},
  \label{eq:return}
\end{equation}
%
where $\gamma \in [0, 1)$ is the \emph{discount factor} controlling the
relative weight of future rewards.

\subsection{Markov Decision Process Formalisation}

The cybersecurity environment in HoneyIQ is formalised as an MDP
$\langle \mathcal{S}, \mathcal{A}, P, R, \gamma \rangle$:
%
\begin{itemize}
  \item $\mathcal{S}$ --- continuous 24-dimensional observation space;
  \item $\mathcal{A}$ --- five discrete honeypot actions;
  \item $P(s' \mid s, a)$ --- transition dynamics determined by the
        attacker's Markov chain (the defender's action does not influence
        the attacker's next state);
  \item $R(s, a, s')$ --- reward function encoding honeypot domain
        knowledge;
  \item $\gamma = 0.99$.
\end{itemize}

\subsection{Q-Learning}

Q-learning estimates the \emph{action-value function} $Q(s, a)$ --- the
expected return from taking action $a$ in state $s$ and then following the
optimal policy:
%
\begin{equation}
  Q^*(s, a) = r + \gamma \cdot \max_{a'} Q^*(s', a')
  \quad \text{(Bellman optimality equation).}
  \label{eq:bellman}
\end{equation}
%
The greedy policy derived from $Q^*$ is $\pi^*(s) = \arg\max_a Q^*(s, a)$.

\section{Deep Q-Networks (DQN)}

Plain Q-learning requires a table indexed by every $(s, a)$ pair, which is
infeasible for continuous or high-dimensional state spaces.
DQN~\citep{mnih2015humanlevel} approximates $Q(s, a; \theta)$ with a neural
network and introduces three stabilising mechanisms.

\subsection{Neural Network Approximator}

HoneyIQ uses a fully connected network:
%
\begin{equation}
  \text{Input}(24)
  \xrightarrow{\text{Linear}} \xrightarrow{\text{LayerNorm}} \xrightarrow{\text{ReLU}}
  \cdots \times 3
  \xrightarrow{\text{Linear}} \mathbb{R}^5.
\end{equation}
%
Hidden layer widths are 256, 128, and 64 respectively.
LayerNorm is preferred over BatchNorm because it stabilises training with
small batch sizes and avoids cross-sample dependencies.
Kaiming uniform initialisation preserves gradient variance under ReLU
activations.

\subsection{Experience Replay}

Transitions $(s, a, r, s', \text{done})$ are stored in a circular replay
buffer of capacity 15\,000.
Mini-batches of 64 are sampled uniformly for each gradient update,
breaking temporal correlations and allowing each transition to be reused
multiple times.

\subsection{Target Network}

A second \emph{target network} with parameters $\theta^-$ provides stable
TD targets:
%
\begin{equation}
  y_i = r_i + \gamma \cdot \max_{a'} Q(s'_i, a'; \theta^-).
  \label{eq:td_target}
\end{equation}
%
$\theta^-$ is hard-copied from the policy network every 150 steps,
preventing the moving-target oscillations that arise when the same network
provides both predictions and targets.

\subsection{Huber Loss and Gradient Clipping}

The TD error is minimised with SmoothL1 (Huber) loss:
%
\begin{equation}
  \mathcal{L}(\delta) =
  \begin{cases}
    \tfrac{1}{2}\delta^2 & |\delta| \le 1 \\
    |\delta| - \tfrac{1}{2} & \text{otherwise,}
  \end{cases}
\end{equation}
%
which is less sensitive to outlier Q-value errors than MSE while remaining
differentiable everywhere.
Gradients are clipped to a maximum $\ell_2$ norm of 10.0 to prevent
divergence.

\subsection{Epsilon-Greedy Exploration}

During training, the agent selects a random action with probability
$\varepsilon$ (exploration) and the greedy action otherwise.
$\varepsilon$ is annealed exponentially from 1.0 to 0.05:
%
\begin{equation}
  \varepsilon \leftarrow \max(\varepsilon_{\min},\; \varepsilon \cdot \varepsilon_{\text{decay}}),
  \quad \varepsilon_{\text{decay}} = 0.997.
\end{equation}

\section{Markov Chains}

An attacker's behaviour is modelled as a Discrete-Time Markov
Chain~\citep{norris1998markov} (DTMC):
%
\begin{equation}
  P(X_{t+1} = j \mid X_t = i,\, X_{t-1}, \ldots) = P(X_{t+1} = j \mid X_t = i).
\end{equation}
%
Two independent chains operate in parallel: a $10 \times 10$ attack-type
transition matrix and a $7 \times 7$ kill-chain-stage transition matrix.
Intent-specific modifier matrices are applied element-wise before
row-renormalisation, skewing the probability mass towards intent-appropriate
transitions without altering the matrix dimensions.

\section{Lockheed Martin Cyber Kill Chain}

The Cyber Kill Chain~\citep{hutchins2011intelligence} models adversary
campaigns as a sequence of seven phases (Table~\ref{tab:killchain}).
HoneyIQ assigns each attack type a \emph{primary stage} and uses stage
weights that increase linearly from 0.10 (Reconnaissance) to 1.00 (Actions
on Objectives), reflecting that later-stage activity is more dangerous.

\begin{table}[ht]
  \centering
  \caption{Cyber Kill Chain stages and associated primary attack types.}
  \label{tab:killchain}
  \begin{tabular}{clll}
    \toprule
    Index & Stage & Weight & Primary Attack Types \\
    \midrule
    0 & Reconnaissance & 0.10 & NORMAL, RECONNAISSANCE \\
    1 & Weaponization  & 0.20 & ANALYSIS \\
    2 & Delivery       & 0.35 & FUZZERS, GENERIC \\
    3 & Exploitation   & 0.55 & EXPLOITS, SHELLCODE \\
    4 & Installation   & 0.70 & BACKDOORS \\
    5 & Command \& Control & 0.85 & WORMS \\
    6 & Actions on Objectives & 1.00 & DOS \\
    \bottomrule
  \end{tabular}
\end{table}

\section{Honeypot Technology}

A honeypot is a decoy system designed to attract, observe, and study
attackers~\citep{spitzner2003honeypot}.
Any connection to a properly deployed honeypot is inherently suspicious
because legitimate users have no reason to interact with it.
HoneyIQ abstracts honeypot responses into five actions:
ALLOW, LOG, TROLL (tarpit), BLOCK, and ALERT.

\section{UNSW-NB15 Dataset}

The UNSW-NB15 benchmark~\citep{moustafa2015unsw} contains network flow
records across nine attack categories and normal traffic, with 49 features
derived from raw packet captures.
HoneyIQ uses 15 of these features, sampling them from parametric
distributions that reproduce per-attack signatures (e.g.\ high
\texttt{ct\_dst\_ltm} for Reconnaissance, extreme \texttt{sload} for DoS).

\section{Random Forests}

A Random Forest~\citep{breiman2001random} is an ensemble of decision trees,
each trained on a bootstrap sample with random feature subsets at each split.
The attack classifier in HoneyIQ uses
\texttt{class\_weight=\textquotedbl{}balanced\textquotedbl{}} to handle
class imbalance and \texttt{StandardScaler} normalisation on the 15-feature
input vector.
Training data is generated entirely from the parametric feature
distributions, avoiding dependence on a real labelled dataset.

\section{Related Work}

Several studies have applied RL to cybersecurity scenarios.
\citet{li2019mad} used multi-agent RL for network defence, demonstrating that
learned policies outperform static rules under adaptive adversaries.
Prior work on DQN-based intrusion response typically treats the attack stream
as a stationary distribution; HoneyIQ extends this by modelling attacker
intent as a latent variable that governs the Markov transition structure,
enabling evaluation of policy robustness across qualitatively distinct
campaigns.
